\usepackage{amsmath,amsfonts,bm}
\usepackage{booktabs} % For formal tables
\usepackage{graphicx}
\usepackage{graphics}
\usepackage{array}
\usepackage{color}
\usepackage{xcolor}
\usepackage{colortbl}
\usepackage{makecell}
\usepackage{multirow}
\usepackage{arydshln}
\usepackage{xspace}
\usepackage{subcaption}
\usepackage{caption}
\usepackage{algpseudocode}
\usepackage[noend,ruled]{algorithm2e}
\usepackage{setspace}
\usepackage{varwidth}
\usepackage{enumitem}
\usepackage{mdframed}
\usepackage{centernot}
\usepackage{textcomp}
\usepackage{soul}
\usepackage{pifont}
\usepackage{wrapfig}
\usepackage{url}
\usepackage{hyperref}
\usepackage[nameinlink]{cleveref}
\usepackage{cuted}
\usepackage{tcolorbox}
\tcbuselibrary{breakable}
\usepackage{lipsum}
\usepackage{listings}
\usepackage{twemojis}
\usepackage[T1]{fontenc}

\usepackage{tikz}
\usetikzlibrary{positioning}

\newcounter{timelineitem}
\newcommand{\timelineentry}[2]{%
  \stepcounter{timelineitem}%
  \node[circle,fill=black,inner sep=2pt] (dot\thetimelineitem) at (0,-\value{timelineitem}) {};
  \node[anchor=west] at (0.7,-\value{timelineitem}) {#1};
  \node[anchor=west] at (5.8,-\value{timelineitem}) {#2};
  \ifnum\value{timelineitem}>1
    \draw[thin] (dot\number\numexpr\value{timelineitem}-1\relax) -- (dot\thetimelineitem);
  \fi
}

\newcommand{\cmark}{\ding{51}}%
\newcommand{\xmark}{\ding{55}}%

\newcommand{\cmtransours}{CMTrans\xspace}
\newcommand{\CC}{Comp-Corp\xspace}

\newcommand{\repostmethodName}{\textsc{RepoST}\xspace}
\newcommand{\repostdatasetName}{\textsc{RepoST}\xspace}
\newcommand{\reposttrainsetName}{\textsc{RepoST-Train}\xspace}
\newcommand{\repostevalsetName}{\textsc{RepoST-Eval}\xspace}

\newcommand{\anchordrours}{Anchor-DR\xspace}

\newcommand{\repostmethodname}{\textsc{Hybrid-Gym}\xspace}

\newcommand{\fencemodelname}{\textsc{FenCE}\xspace}

\newcommand{\plusname}{\textsc{Hybrid-Plus}\xspace}

\newcommand{\proposedWork}{\textbf{\textcolor{treegreen}{proposed work}}}

% circled
\newcommand{\circled}[1]{\textcircled{\raisebox{-0.9pt}{#1}}}



% newly added
\newcommand{\edit}[1]{\textcolor{darkgreen}{#1}}

\newcommand{\TBD}[1]{\textcolor{blue}{[TBD: #1]}}
\newcommand{\yiqing}[1]{\textcolor{blue}{[Yiqing: #1]}}
\newcommand{\daniel}[1]{\textcolor{brown}{[Daniel: #1]}}

\newcommand{\plan}{\textcolor{blue}{[Plan]}}

\newcommand{\ie}{\textit{i.e.},} % that is
\newcommand{\eg}{\textit{e.g.},} % for example
\newcommand{\start}[1]{\vspace{1.5mm}\noindent{{\bf #1}.}}
\DeclareMathOperator*{\argmin}{arg\,min\,}
\DeclareMathOperator*{\argmax}{arg\,max\,}
\newcommand\numberthis{\addtocounter{equation}{1}\tag{\theequation}}


% defining vspace 
\newcommand{\downv}{\vspace{-.0cm}}
\newcommand{\upv}{\vspace{-.0cm}}
% \newcommand{\downve}{\vspace{-.0cm}}


% \usepackage{color, colortbl}
\definecolor{amber}{rgb}{1.0, 0.75, 0.0}
\definecolor{applegreen}{rgb}{0.55, 0.71, 0.0}
\definecolor{treegreen}{rgb}{0.13, 0.54, 0.23}
\definecolor{LightCyan}{rgb}{0.88,1,1}
\definecolor{LightBlue}{RGB}{171, 227, 235}
\newcommand{\bad}{\cellcolor{red!8}}
\newcommand{\worse}{\cellcolor{red!16}}
\newcommand{\good}{\cellcolor{applegreen!10}}
\newcommand{\better}{\cellcolor{applegreen!20}}
\newcommand{\hlcell}{\cellcolor{applegreen!22}}
\newcommand{\hlrow}{\rowcolor{amber!18}}
\newcommand{\dlrow}{\rowcolor{gray!20}}
\newcommand{\sdlrow}{\rowcolor{gray!10}}
\newcommand{\blrow}{\rowcolor{LightBlue!30}}
\newcommand{\dlcell}{\cellcolor{gray!20}}



\definecolor{green2}{HTML}{BFD8B6}
\definecolor{green3}{HTML}{E7F0E5}
\definecolor{greenarrow}{HTML}{1DB100}
\definecolor{red3}{HTML}{C82506}

% defining colours to be used in figures
% \definecolor{gred}{RGB}{219,68,55}
\definecolor{gred}{RGB}{255,102,102}
\definecolor{gblue}{RGB}{51,102,255}
\definecolor{gyellow}{RGB}{244,180,0}
\definecolor{ggreen}{RGB}{15,157,88}
\definecolor{ggrey}{RGB}{115,115,115}
\definecolor{na}{gray}{0.9}

\definecolor{textRed}{RGB}{157,0,23}
\definecolor{textYellow}{RGB}{166,119,54}
\definecolor{textGreen}{RGB}{58,110,38}
\definecolor{textBlue}{RGB}{39,71,156}

% \definecolor{lightRed}{RGB}{255,222,222}
% \definecolor{lightYellow}{RGB}{255,255,193}
% \definecolor{lightGreen}{RGB}{222,255,222}
% \definecolor{lightBlue}{RGB}{221,231,244}

\definecolor{LightYellow}{RGB}{255,250,208}
\definecolor{LightGreen}{RGB}{194,255,192}
\definecolor{LightBlue}{RGB}{187,236,251}
\definecolor{LightPurple}{RGB}{224,223,255}
\definecolor{LightGrey}{RGB}{225,225,225}

\definecolor{OrangeRed}{rgb}{1.0, 0.27, 0.0}
\definecolor{mediumgreen}{RGB}{72, 158, 52}
\definecolor{darkgreen}{rgb}{0.0, 0.42, 0.24}
\definecolor{midnightgreen}{rgb}{0.0, 0.29, 0.33}
\definecolor{aquablue}{RGB}{52, 147, 158}

\definecolor{diagramRed}{RGB}{246,193,193}
\definecolor{diagramPurple}{RGB}{224,224,253}
\definecolor{diagramOrange}{RGB}{244,222,176}
\definecolor{diagramGrey}{RGB}{236,236,236}

\definecolor{Grey}{RGB}{150,150,150}

\newcommand{\ctext}[2]{%
  \begingroup
  \sethlcolor{#1}%
  \hl{ #2 }%
  \endgroup
}

\newcommand{\hlpgen}[2]{\fboxsep1pt \colorbox{ggreen!#2}{\strut #1}} % highlight text with custom saturation. for pgen
\newcommand{\hlcov}[2]{\fboxsep1pt \colorbox{gyellow!#2}{\strut #1}} % highlight text with custom saturation. for cov
\newcommand{\colorR}[1]{\textcolor{gred}{\textbf{#1}}}
\newcommand{\colorG}[1]{\textcolor{ggreen}{\textbf{#1}}}
\newcommand{\colorB}[1]{\textcolor{gblue}{\textbf{#1}}}
\newcommand{\colorGrey}[1]{\textcolor{ggrey}{\textbf{#1}}}

\newcommand{\colorr}[1]{\textcolor{gred}{#1}}
\newcommand{\colorg}[1]{\textcolor{ggreen}{#1}}
\newcommand{\colorb}[1]{\textcolor{gblue}{#1}}

\newcommand{\colordg}[1]{\textcolor{darkgreen}{#1}}
\newcommand{\colorgrey}[1]{\textcolor{Grey}{#1}}